\section{Proofs of the finite-sample results}\label{app:theory-proofs}

The arguments below use fixed design counts. In the IID interpretation,
they apply after legitimate conditioning on the target and reference
positive counts. Independence across audit observations is assumed where
concentration is invoked. No independence between the resulting family
of p-values is required.

\subsection{Fixed-prior risk and the screening boundary}

\begin{proof}[Proof of Proposition~\ref{prop:theory-shrinkage}]
Let $v=p(1-p)$. The posterior mean has bias
$\kappa(\mu-p)/(n+\kappa)$ and variance $nv/(n+\kappa)^2$. Hence
\begin{align*}
 &R_\kappa(p)-R_0(p)\\
 &\quad=\frac{nv+\kappa^2(p-\mu)^2}{(n+\kappa)^2}-\frac vn\\
 &\quad=\frac{\kappa^2(p-\mu)^2+nv-v(n+\kappa)^2/n}{(n+\kappa)^2}\\
 &\quad=\frac{\kappa^2(p-\mu)^2-(2\kappa+\kappa^2/n)v}
 {(n+\kappa)^2},
\end{align*}
which proves Eq.~\eqref{eq:theory-risk}, including $p=0$ or $p=1$.
The strict screening event satisfies
\[
 \widetilde p>t
 \quad\Longleftrightarrow\quad
 X>nt+\kappa(t-\mu).
\]
Since $X$ takes integer values, the smallest count that passes the screen is
$\lfloor nt+\kappa(t-\mu)\rfloor+1$. This proves
Eq.~\eqref{eq:theory-cutoff}. When $t>\mu$, the real threshold increases;
the integer threshold cannot decrease but may remain unchanged. Under
Eq.~\eqref{eq:theory-collapse}, the real threshold is at least $n$, making
the strict event impossible. Conversely, if the inequality fails, $X=n$
passes; thus the condition exactly characterizes collapse of the event
over all possible counts.

If $p=t+\varepsilon$, the shrunk event is $X>np+a$ with
$a=\kappa(t-\mu)-n\varepsilon$. For $a>0$, Hoeffding's inequality gives
\[
 \begin{aligned}
 \Pr_p(A_\kappa)&\le\Pr_p(X-np\ge a)\\
 &\le e^{-2a^2/n}.
 \end{aligned}
\]
For the raw screen,
$A_0^c=\{X/n\le t\}=\{X/n-p\le-\varepsilon\}$, whose probability is at
most $e^{-2n\varepsilon^2}$. These are the bounds in
Eq.~\eqref{eq:theory-screen-bound}. If $\mu,\kappa$ were learned from
independent data, the same argument holds after conditioning on those
data, with the risk interpreted conditionally.
\end{proof}

For the constructive family, ordinary cells have $p=\mu\in(0,1)$, so
Eq.~\eqref{eq:theory-risk} is strictly negative there. The risk difference
for one exceptional cell is finite. Increasing the finite number of
ordinary cells therefore makes the macro-average difference negative,
while Eq.~\eqref{eq:theory-collapse} can still hold for the exceptional
cell. In the numerical example, the ordinary-cell raw and shrunk risks
are $0.009$ and $0.00009$. The exceptional-cell risks are $0.02275$ and
$0.0508525$. Averaging over four ordinary cells and one exceptional cell gives
$0.01175$ and $0.0102425$, respectively. The raw selection probability is
the exact tail
$1-\sum_{k=0}^2 {10\choose k}0.35^k0.65^{10-k}$, while
$nt+\kappa(t-\mu)=2+9=11$ makes the shrunk event empty.

\subsection{HB validity and standardized references}

\begin{proof}[Proof of Proposition~\ref{prop:theory-hb}]
Hoeffding's inequality for the independent, possibly heterogeneous
reference observations gives
\[
 \Pr\left\{q-\widehat q>
       \sqrt{\frac{\log(1/\eta)}{2m}}\right\}\le\eta.
\]
Clipping the endpoint at one preserves coverage because $q\le1$.
Consequently, $E_q=\{q\le U_q\}$ has probability at least $1-\eta$.
Under the null, let $\theta_0=\min\{1,q+\Delta\}$, so
$p\le\theta_0$. On $E_q$, $\theta_0\le\theta_U$. Monotone binomial
coupling implies $T_{\theta_0}(x)\le T_{\theta_U}(x)$ for every count $x$.

For completeness, $T_{\theta_0}(X)$ is super-uniform under any
$p\le\theta_0$. For any $v\in[0,1]$, the set
$\{x:T_{\theta_0}(x)\le v\}$ is empty or an upper integer tail starting at
some $k$. In the latter case,
\[
 \begin{aligned}
 \Pr_p\{T_{\theta_0}(X)\le v\}&=\Pr_p(X\ge k)\\
 &\le T_{\theta_0}(k)\le v.
 \end{aligned}
\]
This includes degenerate binomial parameters.

For $u<\eta$, $\{P_{\mathrm{HB}}\le u\}$ is empty. For
$\eta\le u<1$, clipping at one does not change this event, and
\[
 \begin{aligned}
 &\{P_{\mathrm{HB}}\le u\}\\
 &\quad\subseteq E_q^c\ \cup\
 \{T_{\theta_0}(X)\le u-\eta\}.
 \end{aligned}
\]
The union bound gives a probability of at most $\eta+(u-\eta)=u$.
The case $u=1$ is immediate. A p-value that is identically one for an unsupported
comparison is also super-uniform. This proves the proposition. The
argument does not condition the binomial law of $X$ on the random event
$E_q$. Although target/reference independence holds in the stated audit
design, this union-bound step itself requires only the binomial target
marginal and the valid reference bound.
\end{proof}

For the standardized-reference extension, write observations from cell
$h$ as $B_{hi}$, $i=1,\ldots,n_h$. Then
$\widehat q=\sum_h\sum_i(\nu_h/n_h)B_{hi}$. The sum of squared ranges is
$\sum_h n_h(\nu_h/n_h)^2=\sum_h\nu_h^2/n_h$. The weighted Hoeffding inequality therefore
gives the stated upper radius for $q=\sum_h\nu_hp_h$. The preceding proof
then applies unchanged. This argument requires $n_h>0$ whenever
$\nu_h>0$; silently discarding unsupported reference cells would change
the target and invalidate the argument.

For Eq.~\eqref{eq:theory-hoeffding}, express $d$ as a sum of independent
target observations with coefficient $1/n$ and reference observations
with coefficient $-1/m$. The sum of squared ranges is $V=1/n+1/m$.
If the difference between the average means satisfies $\delta\le\Delta$, then, for
$x>0$,
\[
 \begin{aligned}
 \Pr(d\ge\Delta+x)&\le\Pr(d-\delta\ge x)\\
 &\le e^{-2x^2/V}.
 \end{aligned}
\]
For $0<u<1$, the event $P_{\mathrm H}\le u$ implies
$d\ge\Delta+\sqrt{(V/2)\log(1/u)}$ and has probability at most $u$.
At $u=0$ it is empty, and at $u=1$ the assertion is automatic. This proves
super-uniformity for heterogeneous rates on both sides.

\subsection{Weighted Holm, the power envelope, and split cost}

\begin{proof}[Proof of Proposition~\ref{prop:theory-holm}]
Condition on the discovery data and the legitimate confirmation design. Weights
are then fixed and strictly positive. Let $I_0$ be the true-null set. If
it is empty, FWER is zero. Otherwise, let $W_0=\sum_{g\in I_0}w_g>0$.
If the procedure rejects any true null, consider the first such
hypothesis $g$ in the ordered sequence. All true nulls remain at its
testing stage. Its remaining-weight denominator $R_g$ therefore satisfies
$R_g\ge W_0$, and rejection implies
\[
 P_g\le\alpha w_g/R_g\le\alpha w_g/W_0.
\]
Taking a union over possible first true nulls and using conditional
super-uniformity gives
\[
 \begin{aligned}
 &\Pr(\text{any false flag}\mid D,\text{design})\\
 &\quad\le\sum_{g\in I_0}\frac{\alpha w_g}{W_0}=\alpha.
 \end{aligned}
\]
Each summand uses a threshold at most $\alpha\le1$. Averaging over the
conditioning variables proves strong unconditional control. The argument does not use joint
independence of p-values.

For containment, suppose $P_g/w_g\le\alpha$. All preceding ratios are at
most $\alpha$. At any such stage the remaining-weight sum is at most
one, so the Holm ratio threshold is at least $\alpha$. The procedure
reaches and rejects $g$, proving weighted-Bonferroni containment.
\end{proof}

To prove Eq.~\eqref{eq:theory-power-envelope}, again condition on
discovery and the fixed design. With $P_{\mathrm H}$, the event
$d_g\ge\Delta+b_g(w_g)$ implies $P_g\le\alpha w_g$, hence a weighted
Holm rejection. Under $\delta_g=\Delta+\varepsilon_g$ and
$\varepsilon_g>b_g(w_g)$, a lower-tail Hoeffding bound gives
\[
 \begin{aligned}
 &\Pr\{d_g<\Delta+b_g(w_g)\}\\
 &\quad\le\exp\left\{-\frac{2[\varepsilon_g-b_g(w_g)]^2}{V_g}\right\}.
 \end{aligned}
\]
Taking the complementary probability proves the envelope. Since
$w_g\ge\rho/G$, using the larger boundary $b_g(\rho/G)$ gives its
score-independent version whenever that bound is nontrivial. Neither
calculation asserts that HB has the same boundary or that a zero lower
bound implies zero actual power.

\begin{proof}[Proof of Corollary~\ref{cor:theory-split}]
Exactly proportional allocation gives $n_C=cn_A$, $m_C=cm_A$, and hence
$V_C=V_A/c$. The split and full-data sufficient boundaries satisfy
\[
 \begin{aligned}
 b_C(w_g)^2&=\frac{V_A}{2c}\log\frac1{\alpha w_g},\\
 b_A(1/G)^2&=\frac{V_A}{2}\log\frac G\alpha.
 \end{aligned}
\]
Both boundaries are nonnegative. Comparing their squares and cancelling
$V_A/2$ yields
\[
 \begin{aligned}
 &b_C(w_g)<b_A(1/G)\\
 &\quad\Longleftrightarrow\quad
 \log\frac1{\alpha w_g}<c\log\frac G\alpha\\
 &\quad\Longleftrightarrow\quad
 w_g>\alpha^{c-1}G^{-c}.
 \end{aligned}
\]
This proves the asserted equivalence solely for the specified sufficient
boundaries.
\end{proof}
